\documentclass{ctexart}

\usepackage{amsmath,amsthm,extarrows,color,enumitem,framed,amscd}
\usepackage[a4paper, left = 1cm, right = 1cm]{geometry}

\usepackage{tikz}
\usetikzlibrary{arrows.meta}
\usepackage{tikz-cd}

\usepackage[texMathDollars]{markdown}

% 取消首行缩进
\ctexset{autoindent=false}   % 先关闭 ctex 的自动调整功能
\setlength{\parindent}{0pt}  % 再把段首缩进设为 0

\usepackage{unicode-math}
%\setmathfont{STIX Two Math}
%\setmathfont{TeX Gyre Termes Math}
%\setmathfont{Libertinus Math}
%\setmathfont{Fira Math}

%%%%%%%% FONT TIMESAVES %%%%%%%
\newcommand{\bT}{\mathbb{T}}
\newcommand{\bN}{\mathbb{N}}
\newcommand{\bI}{\mathbb{I}}

\newcommand{\CC}{\mathcal{C}}

%%%%%%%% CATEGORY THEORY %%%%%%%
\newcommand{\id}{\mathrm{id}}
\DeclareMathOperator{\Hom}{Hom}

\newcommand{\cat}[1]{\mathrm{#1}}
\newcommand{\Set}{\cat{Set}}

\DeclareFontFamily{U}{dmjhira}{}
\DeclareFontShape{U}{dmjhira}{m}{n}{ <-> dmjhira }{}

\DeclareRobustCommand{\yo}{\text{\usefont{U}{dmjhira}{m}{n}\symbol{"48}}}

\DeclareMathOperator*{\colim}{colim}

\newcommand{\inl}{\operatorname{inl}}
\newcommand{\inr}{\operatorname{inr}}

%%%%%%%%% GEOMTERY %%%%%%%%%%%%%%
\newcommand{\Spec}{\operatorname{Spec}}
\newcommand{\Zar}{\mathrm{Zar}}

%%%%%%%%% TYPE THEORY %%%%%%%%%%%
\newcommand{\OO}{\bigcirc}
\newcommand{\UU}{\mathcal{U}}
\newcommand{\isEquiv}{\bibopenparen{isEquiv}}
\newcommand{\hlevel}{\mathrm{hlevel}}
\newcommand{\isContr}{\operatorname{isContr}}
\newcommand{\isProp}{\operatorname{isProp}}
\newcommand{\Prop}{\mathrm{Prop}}
\newcommand{\refl}{\mathrm{refl}}
\newcommand{\isModal}{\operatorname{isModal}}
\newcommand{\ap}{\operatorname{ap}}
\newcommand{\fib}{\operatorname{fib}}
\newcommand{\im}{\operatorname{im}}
\newcommand{\tot}{\mathrm{tot}}
\newcommand{\pt}{\mathrm{pt}}

\begin{document}
\begin{markdown}
设 $\kappa$ 为一个**正则基数**（regular cardinal），**$\kappa$-滤过余极限**（$\kappa$-filtered colimit）是指定义在 **$\kappa$-滤过范畴**（$\kappa$-filtered category）上的图表的余极限（colimit）。

**1. $\kappa$-滤过范畴的精确定义**

一个小范畴（或局部小范畴） $\mathcal{I}$ 称为 **$\kappa$-滤过范畴**，如果它满足以下条件之一（两种经典等价定义）：

**定义 A（元素与态射有界性）**：

* **对象有上界**：对于任意势严格小于 $\kappa$ 的对象集合 $S \subseteq \text{Ob}(\mathcal{I})$（即 $\vert{}S\vert{} < \kappa$），存在一个对象 $j \in \mathcal{I}$ 及一族态射 $\{f_i: i \to j\}_{i \in S}$。
* **平行态射可均衡**：对于任意两个对象 $i, j \in \mathcal{I}$ 以及任意势严格小于 $\kappa$ 的平行态射集合 $R \subseteq \text{Hom}_{\mathcal{I}}(i, j)$（即 $\vert{}R\vert{} < \kappa$），存在对象 $k \in \mathcal{I}$ 及态射 $g: j \to k$，使得对任意 $f_1, f_2 \in R$，均有：

$$g \circ f_1 = g \circ f_2$$



**定义 B（范畴图表余锥性）**：
对于任意态射集势严格小于 $\kappa$ 的小范畴 $\mathcal{D}$（即 $\vert{}\text{Mor}(\mathcal{D})\vert{} < \kappa$），任何函子 $F: \mathcal{D} \to \mathcal{I}$ 在 $\mathcal{I}$ 中都存在一个余锥（cocone）。

---

**2. $\kappa$-滤过余极限的定义**

给定范畴 $\mathcal{C}$ 与函子 $D: \mathcal{I} \to \mathcal{C}$：

* 当索引范畴 $\mathcal{I}$ 为 $\kappa$-滤过范畴时，图表 $D$ 的余极限 $\varinjlim_{\mathcal{I}} D$（或记作 $\text{colim}_{i \in \mathcal{I}} D(i)$）即被称为 **$\kappa$-滤过余极限**。

---

**3. 特例与经典表达**

* **$\aleph_0$-滤过余极限**：当 $\kappa = \aleph_0$ 时，$\aleph_0$-滤过范畴即为标准的**滤过范畴**（filtered category），其余极限即为传统的**滤过余极限**（如环的定向极限、模的有向并等）。
* **集合范畴 $\mathbf{Set}$ 中的具体构造**：若 $D: \mathcal{I} \to \mathbf{Set}$ 为 $\kappa$-滤过图表，其余极限由不相交并的等价类给出：

$$\varinjlim_{i \in \mathcal{I}} D(i) = \left( \bigsqcup_{i \in \text{Ob}(\mathcal{I})} D(i) \right) \Big/ \sim$$



其中等价关系 $\sim$ 定义为：对 $x \in D(i)$ 与 $y \in D(j)$，$x \sim y$ 当且仅当存在 $k \in \mathcal{I}$ 以及态射 $f: i \to k, g: j \to k$，使得在 $D(k)$ 中满足：

$$D(f)(x) = D(g)(y)$$
\end{markdown}
\end{document}
